\documentclass[12pt]{article}

\usepackage{sbc-template}
\usepackage{graphicx,url}
\usepackage[utf8]{inputenc}
\usepackage[brazil]{babel}
\usepackage{cite}
\usepackage{booktabs}
\usepackage{multirow}
\usepackage{hyperref}

\sloppy

\title{Mapeamento de Biomarcadores e Estimadores de Estados Mentais em Interfaces Cérebro-Computador Passivas: Uma Revisão de Escopo}

\author{Felippe Veloso Marinho\inst{1}\thanks{ORCID: \url{https://orcid.org/0009-0003-6910-9143}},
        Coautor(a) do Trabalho\inst{1}}

\address{Escola de Engenharia\\
  Universidade Federal de Minas Gerais (UFMG) -- Belo Horizonte, MG -- Brasil\\
  \email{felippevm@ufmg.br}}

\begin{document}

\maketitle

\begin{abstract}
Passive brain--computer interfaces (pBCIs) decode covert cognitive and affective states to enable neuroadaptive systems without requiring explicit user commands. This scoping review maps neurophysiological biomarkers, machine learning pipelines, and methodological validity criteria reported in studies that construct or validate quantitative EEG-based mental-state estimators within passive BCI frameworks. Following PRISMA-ScR guidelines, we searched OpenAlex and PubMed (January 2010--March 2026), screened records with a priori inclusion rules and a structured relevance score, and extracted features across temporal (ERPs), spectral (theta, alpha, mu, beta), spatial, and connectivity domains. Emphasis is placed on cross-validation practices that account for temporal dependence, calibration and cross-task generalization, and the neurophysiological interpretability of linear and nonlinear decoders. Of 82 unique records screened, 13 were pre-selected for full-text synthesis; preliminary mapping indicates dominant spectral ratios and engagement indices, heterogeneous validation rigor, and limited reporting of prefrontal channels and pedagogical applications.
\end{abstract}

\begin{resumo}
Interfaces cérebro-computador passivas (pBCIs) decodificam estados cognitivos e afetivos implícitos para viabilizar sistemas neuroadaptativos sem exigir comandos voluntários explícitos. Esta revisão de escopo mapeia biomarcadores neurofisiológicos, \textit{pipelines} de aprendizado de máquina e critérios de validade metodológica em estudos que constroem ou validam estimadores quantitativos de estados mentais baseados em EEG no paradigma pBCI. Seguindo PRISMA-ScR, buscamos OpenAlex e PubMed (janeiro/2010--março/2026), triamos registros com critérios \textit{a priori} e pontuação estruturada de relevância, e extraímos \textit{features} nos domínios temporal (ERPs), espectral (teta, alfa, mu, beta), espacial e de conectividade. Destacam-se práticas de validação cruzada sensíveis à dependência temporal, generalização entre tarefas e sessões, e interpretabilidade neurofisiológica de decodificadores lineares e não lineares. Dos 82 registros únicos triados, 13 foram pré-selecionados para síntese em texto integral; o mapeamento preliminar indica predominância de razões espectrais e índices de engajamento, rigor de validação heterogêneo e escassa reportagem de canais frontopolar e aplicações pedagógicas.
\end{resumo}

\section{Introdução}

O monitoramento contínuo de estados cognitivos e afetivos por meio de interfaces cérebro-computador passivas (\textit{passive brain--computer interfaces} --- pBCIs) consolidou-se como eixo central da neuroergonomia e da interação humano-máquina adaptativa~\cite{Arico2017,Zander2016}. Diferentemente dos paradigmas ativos ou reativos, os sistemas passivos inferem informação a partir da atividade cerebral espontânea ou induzida pela tarefa principal, sem sobrecarregar o usuário com comandos voluntários explícitos~\cite{Zander2016}. Esse paradigma viabiliza aplicações em sistemas de alta complexidade e criticidade de segurança --- como aviação, controle de tráfego aéreo (ATM), condução veicular e ambientes clínicos --- nos quais a interface deve adaptar comportamento e funcionalidade ao estado mental real do operador~\cite{Arico2017,Arico2016ATM}.

Aricò \textit{et al.}~\cite{Arico2017}, em \textit{mini-revisão} sobre pBCI em ambientes operacionais, sintetizaram quatro eixos metodológicos recorrentes: (i)~avaliação de estados mentais por \textit{neurometrics} (EEG e, em menor escala, fNIRS, fMRI, MEG e biosinais periféricos); (ii)~tecnologias pBCI para melhoria da interação humano-máquina em \textit{closed-loop}; (iii)~identificação de \textit{features} eletrofisiológicas (oscilações espectrais, ERPs, conectividade); e (iv)~mensuração de estados mentais por aprendizado de máquina, com recomendações explícitas sobre modelos lineares, kernels, regularização e calibração. Estados como estresse, fadiga mental, sonolência, sobrecarga de trabalho (\textit{mental workload}), frustração e engajamento/atencão aparecem como alvos recorrentes; componentes temporais como P300 (recursos atencionais) e ERN (detecção de erro) exemplificam biomarcadores com potencial operacional~\cite{Arico2017,Blankertz2011}.

Contudo, a translação laboratório--campo permanece limitada. Baixa relação sinal-ruído (SNR), artefatos fisiológicos e instrumentais, não estacionariedade do EEG, calibração dependente de tarefa e sessão, e riscos de sobreajuste (\textit{overfitting}) --- frequentemente agravados por validação cruzada que ignora dependência temporal entre janelas adjacentes~\cite{Varoquaux2017} --- comprometem confiabilidade e reprodutibilidade. Adicionalmente, interpretar vetores de peso de classificadores lineares como mapas de ativação neural exige distinção cuidadosa entre modelos diretos, inversos e mapas de sensibilidade~\cite{Haufe2014}; confundir essas quantidades pode levar a conclusões neurofisiológicas espúrias.

Revisões sistemáticas recentes em domínios adjacentes --- como a integração BCI--FES para reabilitação motora~\cite{Faria2026} --- demonstraram que, mesmo em corpora maduros, a reportagem de canais EEG, bandas espectrais e algoritmos de classificação permanece incompleta (respectivamente 90\%, 90\% e 94\% nos 82 estudos incluídos), com predominância de ritmo mu (8--13~Hz), beta (13--30~Hz) e derivações sensorimotoras C3/C4/Cz, enquanto eletrodos frontopolar (Fp1/Fp2) aparecem em apenas 10\% dos trabalhos. Esses achados motivam uma revisão de escopo focada especificamente em \textbf{biomarcadores quantitativos} e \textbf{validação metodológica} no paradigma \textbf{passivo}, incluindo contextos pedagógicos, psicológicos e de reabilitação cognitiva, nos quais a literatura permanece fragmentada.

Este trabalho apresenta uma revisão de escopo conduzida conforme PRISMA-ScR~\cite{Tricco2018}, visando:
\begin{enumerate}
    \item mapear biomarcadores e estimadores de estados mentais nos domínios temporal, espectral, espacial e de conectividade em estudos pBCI;
    \item categorizar arquiteturas de aprendizado de máquina (lineares, baseadas em kernel, Riemannianas e \textit{deep learning}) e procedimentos de construção/validação reportados;
    \item analisar rigor metodológico na validação cruzada, calibração, generalização entre tarefas/sessões e interpretabilidade neurofisiológica dos decodificadores;
    \item identificar lacunas de reportagem --- especialmente canais, razões espectrais, \textit{ground truth} e aplicações pedagógicas --- para orientar diretrizes futuras de pesquisa e translação interdisciplinar.
\end{enumerate}

\section{Métodos}

\subsection{Desenho do estudo}

Esta revisão de escopo foi conduzida conforme a extensão PRISMA para \textit{scoping reviews} (PRISMA-ScR)~\cite{Tricco2018}, com protocolo, critérios de elegibilidade e estratégia de busca definidos \textit{a priori}, seguindo o arcabouço População--Conceito--Contexto (PCC)~\cite{Arksey2005,Tricco2018}:
\begin{itemize}
    \item \textbf{População (P):} indivíduos monitorados quanto a variações de estados cognitivos, atencionais, afetivos ou de engajamento em contextos operacionais, clínicos, psicológicos, pedagógicos ou de reabilitação;
    \item \textbf{Conceito (C):} construção, validação ou uso de biomarcadores neurofisiológicos quantitativos (predominantemente EEG) acoplados a modelos preditivos/decodificadores em sistemas pBCI;
    \item \textbf{Contexto (C):} aplicações ecológicas ou semi-ecológicas, neuroergonomia, pedagogia/aprendizagem, psicologia cognitiva, reabilitação e interação humano-computador adaptativa (\textit{closed-loop}).
\end{itemize}

A unidade de análise é a publicação individual (\textit{full paper}). O escopo complementa revisões prévias do grupo~\cite{Faria2026} ao focar explicitamente no paradigma passivo e na cadeia biomarcador--validação--interpretabilidade, em vez de integração BCI--FES ou controle motor explícito.

\subsection{Estratégia de busca}

Duas bases eletrônicas foram consultadas até 20 de março de 2026: \textbf{OpenAlex}, pela cobertura bibliográfica ampla e interdisciplinar; e \textbf{PubMed}, pela cobertura clínica e biomédica. A seleção de bases seguiu critério de complementaridade semelhante ao adotado em revisão PRISMA-ScR anterior do grupo~\cite{Faria2026}, adaptada ao domínio pBCI e biomarcadores.

A string de busca combinou quatro blocos booleanos (\texttt{AND} entre blocos, \texttt{OR} intra-bloco), derivados dos eixos PCC e dos termos recorrentes na literatura de referência~\cite{Arico2017,Lotte2018}:
\begin{enumerate}
    \item \textbf{pBCI:} \texttt{``passive BCI''}, \texttt{``passive brain-computer''}, \texttt{pBCI}, \texttt{``implicit BCI''};
    \item \textbf{biomarcador quantitativo:} \texttt{biomarker}, \texttt{spectral}, \texttt{``engagement index''}, \texttt{``EEG feature''}, \texttt{quantitative};
    \item \textbf{validação:} \texttt{validation}, \texttt{reliability}, \texttt{validity}, \texttt{``cross-validation''};
    \item \textbf{domínio de aplicação:} \texttt{education}, \texttt{learning}, \texttt{pedagog}, \texttt{classroom}, \texttt{rehabilitation}, \texttt{psychology}, \texttt{``human-computer''}.
\end{enumerate}

Não foi imposta restrição de idioma. O período de elegibilidade foi janeiro de 2010 a 20 de março de 2026, alinhado ao critério temporal I4 empregado em~\cite{Faria2026}. A busca foi executada programaticamente via APIs públicas (OpenAlex REST API; NCBI E-utilities para PubMed), com filtro de data aplicado na consulta e verificação pós-busca pelo ano de publicação. Os registros identificados foram exportados para biblioteca de referências (Zotero) e planilha de triagem. A string completa e \textit{logs} de execução estão disponíveis mediante solicitação aos autores.

\subsection{Critérios de inclusão e exclusão}

Estudos foram incluídos na pré-seleção quando, no título ou resumo, atendiam \textbf{simultaneamente} a:
\begin{description}
    \item[I1] menção explícita a pBCI ou equivalente implícito (\textit{passive/implicit brain--computer interface});
    \item[I2] reporte de biomarcador neurofisiológico quantitativo \textbf{ou} procedimento de validação/construção de estimador (p.ex., confiabilidade, validade de critério, validação cruzada, \textit{ground truth});
    \item[I3] pontuação de relevância $\geq 5$ na triagem estruturada (Seção~\ref{sec:triagem});
    \item[I4] publicação entre janeiro/2010 e 20/março/2026 como artigo completo.
\end{description}

Estudos foram excluídos quando:
\begin{description}
    \item[E1] não mencionavam pBCI ou equivalente no título/resumo;
    \item[E2] não reportavam biomarcador quantitativo nem procedimento de validação;
    \item[E3] tratavam predominantemente de BCI ativo, \textit{motor imagery}, prótese/exoesqueleto/\textit{speller}, epilepsia/crise como foco exclusivo, ou modalidades exclusivamente fNIRS/fMRI/MRI sem EEG;
    \item[E4] apresentavam pontuação abaixo do limiar de relevância;
    \item[E5] eram revisões, editoriais, protocolos sem dados, simulações exclusivas ou resumos de conferência sem texto integral.
\end{description}

A confirmação final de elegibilidade requer leitura integral do texto (Seção~\ref{sec:triagem}).

\subsection{Triagem, deduplicação e extração de dados}
\label{sec:triagem}

\textbf{Deduplicação.} Registros brutos das bases foram deduplicados em duas etapas: (i)~por DOI normalizado; (ii)~por similaridade de título ($\geq 90\%$, \textit{fuzzy matching}), preservando o registro mais completo.

\textbf{Triagem por título e resumo.} Adotou-se triagem em duas camadas. Primeiro, regras determinísticas de inclusão/exclusão (I1--I4, E1--E5). Em seguida, um \textbf{dicionário ponderado de termos} --- calibrado a partir dos eixos metodológicos de Aricò \textit{et al.}~\cite{Arico2017} e dos domínios PCC --- atribuiu pontos a padrões textuais no título e resumo concatenados:
\begin{itemize}
    \item \textit{Passive BCI} (obrigatório): p.ex., \texttt{passive BCI}, \texttt{pBCI}, \texttt{implicit BCI};
    \item \textit{Biomarcador}: p.ex., razão beta/alfa, potência espectral, índice de engajamento, \textit{EEG feature};
    \item \textit{Validação}: p.ex., \textit{reliability}, \textit{cross-validation}, \textit{ground truth}, validade ecológica;
    \item \textit{Domínio}: pedagogia (peso maior), reabilitação, psicologia, IHC, neuroergonomia;
    \item \textit{Penalidades de exclusão}: p.ex., \textit{motor imagery}, \textit{active BCI}, modalidade exclusivamente fMRI/fNIRS.
\end{itemize}

Registros com pontuação total $\geq 5$ e sem penalidade determinante ($\leq -6$) foram classificados como \textbf{INCLUÍDO (Pré-seleção)}; os demais, \textbf{EXCLUÍDO}, com motivo registrado. A triagem automatizada funciona como filtro de sensibilidade; a decisão definitiva cabe à avaliação do texto integral por pelo menos um revisor, com citação \textit{verbatim} para campos sujeitos a interpretação --- prática alinhada à mitigação de viés de revisor único adotada em~\cite{Faria2026}.

\textbf{Extração de dados.} Para estudos pré-incluídos, formulário padronizado com vocabulário fechado registra: identificação; população e desenho; modalidade e canais EEG; domínio de \textit{features} (temporal/espectral/espacial/conectividade); algoritmo de decodificação; procedimentos de construção do biomarcador (normalização, janela temporal, índice composto); validação (tipo de CV, confiabilidade, AUC/ROC); \textit{ground truth}; domínio de aplicação; aspectos de IHC (\textit{closed-loop}, calibração, interpretabilidade); e transferibilidade entre contextos (p.ex., reabilitação $\rightarrow$ pedagogia). Campos não reportados explicitamente são codificados ``NR'' (\textit{not reported}); inferência não é permitida.

\textbf{Prioridade pedagógica.} Estudos pré-incluídos recebem classificação auxiliar: \textbf{Alta} (pedagogia explícita), \textbf{Média} (psicologia/IHC/aprendizagem genérica) ou \textbf{Baixa} (reabilitação, neuroergonomia ou domínio indeterminado), para orientar síntese narrativa sem substituir critérios de inclusão.

\subsection{Critérios de classificação analítica}

Para reduzir subjetividade na síntese, campos categóricos seguem regras explícitas, adaptadas de~\cite{Faria2026}:
\begin{itemize}
    \item \textbf{Domínio de \textit{feature}:} temporal (ERP/MRCP), espectral (bandas ou razões), espacial (CSP, filtros espaciais), conectividade (coerência, PLV, grafo), híbrido ou NR;
    \item \textbf{Família de classificador:} linear (LDA, regressão), kernel (SVM), Riemanniano, \textit{deep learning}, limiar/heurístico ou NR;
    \item \textbf{Rigor de validação:} adequado (CV sensível a dependência temporal ou validação externa), parcial (CV padrão sem controle temporal), insuficiente (treino/teste único ou NR);
    \item \textbf{Interpretabilidade:} mapa direto, transformação de pesos (p.ex.,~\cite{Haufe2014}), não abordada ou NR.
\end{itemize}

\section{Resultados}

\subsection{Seleção de estudos}

As buscas identificaram 86 registros brutos (OpenAlex: 80; PubMed: 6). Após deduplicação, restaram 82 registros únicos triados por título e resumo; 69 foram excluídos e 13 pré-selecionados para leitura integral e extração completa (Fig.~\ref{fig:prisma}). Os motivos de exclusão na triagem automatizada incluíram ausência de pBCI explícito, ausência simultânea de biomarcador e validação, escopo incompatível (BCI ativo/\textit{motor imagery}) e pontuação abaixo do limiar.

\begin{figure}[ht]
  \centering
  \fbox{\parbox{0.92\linewidth}{\centering\vspace{2.2cm}
  \textbf{Figura 1} --- Fluxograma PRISMA-ScR (rascunho).\\
  Identificados: 86 \ $\rightarrow$ Únicos: 82 \ $\rightarrow$ Pré-incluídos: 13\\
  \textit{[Inserir figura final]} \vspace{2.2cm}}}
  \caption{Fluxograma de seleção de estudos (PRISMA-ScR).}
  \label{fig:prisma}
\end{figure}

\subsection{Características preliminares do corpus pré-selecionado}

A Tabela~\ref{tab:resumo} resume distribuições preliminares entre os 13 estudos pré-incluídos. A leitura integral confirmará ou revisará cada classificação.

\begin{table*}[ht]
  \centering
  \caption{Características preliminares dos 13 estudos pré-selecionados (triagem título/resumo).}
  \label{tab:resumo}
  \small
  \begin{tabular}{@{}lrr@{}}
    \toprule
    \textbf{Característica} & \textbf{Categoria} & \textbf{$n$ (\%)} \\
    \midrule
    \multirow{3}{*}{Prioridade pedagógica} & Alta & 8 (62) \\
     & Média & 1 (8) \\
     & Baixa & 4 (31) \\
    \midrule
    \multirow{5}{*}{Domínio principal (triagem)} & Pedagogia & 8 (62) \\
     & Psicologia & 1 (8) \\
     & Reabilitação & 1 (8) \\
     & Neuroergonomia & 1 (8) \\
     & Indeterminado & 2 (15) \\
    \bottomrule
  \end{tabular}
\end{table*}

\subsection{Síntese individual dos estudos}

A Tabela~\ref{tab:estudos} apresenta o esquema de extração por estudo, alinhado à matriz PRISMA do projeto. Linhas completas serão preenchidas após leitura integral; campos seguem vocabulário fechado (NR quando não reportado).

\begin{table*}[ht]
  \centering
  \caption{Esquema de síntese individual dos estudos incluídos (a preencher após leitura integral).}
  \label{tab:estudos}
  \scriptsize
  \begin{tabular}{@{}p{2.1cm}p{3.2cm}p{1.2cm}p{1.5cm}p{2.0cm}p{2.0cm}p{1.8cm}p{1.8cm}p{1.6cm}@{}}
    \toprule
    \textbf{Autor (ano)} & \textbf{Título (abrev.)} & \textbf{$n$} & \textbf{Canais EEG} & \textbf{Biomarcador} & \textbf{Validação} & \textbf{Classificador} & \textbf{Ground truth} & \textbf{Domínio} \\
    \midrule
    Aricò (2017) & pBCI operacional: mini-revisão & NR & NR & Espectral, ERP, conectividade & Revisão narrativa & Linear, kernel & NR & Neuroergonomia \\
    -- & Engajamento em MOOC & NR & NR & Índice de engajamento & NR & NR & NR & Pedagogia \\
    -- & CV e métricas pBCI & NR & NR & NR & Cross-validation & NR & NR & Pedagogia \\
    -- & Sonolência ao volante (razões espectrais) & NR & NR & $R_1$--$R_4$ espectrais & NR & NR & NR & Pedagogia \\
    \multicolumn{9}{c}{\textit{Demais linhas: preencher a partir de \texttt{Matriz\_Mapeamento\_PRISMA.xlsx}}} \\
    \bottomrule
  \end{tabular}
\end{table*}

\textbf{Campos recomendados para a tabela final} (espelhando CBEB~\cite{Faria2026} e a matriz do projeto):
\begin{enumerate}
    \item identificação (autor, ano, DOI);
    \item desenho, $n$ participantes, população;
    \item paradigma pBCI e contexto ecológico;
    \item canais EEG e bandas/\textit{features};
    \item construção do biomarcador (normalização, janela, índice composto);
    \item algoritmo de decodificação e calibração;
    \item validação (tipo de CV, métricas, confiabilidade);
    \item \textit{ground truth} (desempenho, autorrelato, outcome clínico/pedagógico);
    \item domínio, aplicação pedagógica, aspectos IHC e transferibilidade.
\end{enumerate}

\section{Discussão (rascunho)}

\subsection{Síntese narrativa esperada}

Com base nas anotações preliminares e na literatura de referência~\cite{Arico2017,Haufe2014,Lotte2018}, espera-se que a síntese final organize achados nos eixos:
\begin{itemize}
    \item \textbf{Neurometrics e domínios de \textit{feature}:} predominância espectral (teta, alfa, beta, razões como $\alpha/\beta$, $(\alpha+\theta)/\beta$) frente a ERPs (P300, ERN) e conectividade; comparação com a concentração sensorimotora (C3/C4/Cz) observada em revisões de BCI motor~\cite{Faria2026};
    \item \textbf{Aprendizado de máquina:} trade-off viés--variância; adequação de modelos lineares regularizados em calibração escassa versus kernels e \textit{deep learning} em amostras maiores~\cite{Arico2017,Lotte2018};
    \item \textbf{Validação e calibração:} generalização entre tarefas (p.ex., \textit{n-back} vs.\ condução), deriva intersessão e necessidade de readaptação online~\cite{Arico2016ATM,Varoquaux2017};
    \item \textbf{Interpretabilidade:} distinção entre pesos discriminativos e padrões de ativação forward~\cite{Haufe2014};
    \item \textbf{Lacunas:} reportagem de Fp1/Fp2, validação ecológica em pedagogia, e \textit{pipelines} reprodutíveis ponta a ponta.
\end{itemize}

\subsection{Trabalhos futuros}

A consolidação desta revisão em manual metodológico --- com \textit{pipelines} exemplares por domínio (espectral, temporal, espacial/Riemanniano) e guia de escolha de classificador condicionado ao volume de calibração e ao SNR --- permanece como extensão natural, condicionada à maturação empírica dos estudos incluídos.

\section*{Agradecimentos}

% Inserir financiamento (CAPES, FAPEMIG, etc.) conforme CBEB.

\bibliographystyle{sbc}
\bibliography{referencias}

\end{document}
